% !TeX program = platex
% Encoding: UTF-8
\PassOptionsToPackage{dvipdfmx}{color}
\documentclass[submit,techrep]{ipsj}
\usepackage[dvipdfmx]{graphicx}
\usepackage{url}
\usepackage{array}
\setcounter{year}{2026}
\providecommand{\newblock}{\hskip .11em plus .33em minus .07em}
\renewcommand{\doi}[1]{{\normalfont DOI: \urlstyle{same}\url{#1}}}
\begin{document}
\title{公開WebAssembly採掘データセットにおける難読化とハッシュ計算性能の実証分析}
\etitle{An Empirical Analysis of Obfuscation and Hashing Performance in Public WebAssembly Mining Datasets}
\affiliate{AFF}{金沢工業大学大学院\\Graduate School, Kanazawa Institute of Technology\\指導教員：向井宏明}
\author{貝沼英亮}{貝沼英亮}{AFF}
\begin{abstract}
公開WebAssembly採掘データセットの難読化について，変換の適用範囲と，正しいハッシュ計算を実行した際の性能を対応付けて評価する．既存研究の対象と測定条件を照合し，Proof-of-Theftの配布10標本，SWAMPED入力1標本と公開変換コードから生成した2標本を，公表テストベクトルによる出力照合と反復速度測定の対象とした．xmr-wasmでは30秒以上・3反復の速度中央値が配布基準9.11 H/s，平坦化版6.18 H/sであり，命令数の増加量と速度低下の順位は置換版との間で一致しなかった．また，命令列で確認した変換の一部は今回のv0入力では通らず，静的な増加量だけでは実行負荷を説明できない．SWAMPEDの同一組立条件の自生成対では，nop挿入・整数定数分割の命令数比が1.097・1.034，速度中央値比が0.960・0.973で，反復範囲が重なった．配布設定の不一致，最適化条件の違い，同一標本の再収録も確認し，手法の因果効果や独立した実装での再現と区別した．追加実測は外部プールを用いないローカル環境で行った．本稿は手法の一般的な説明ではなく，公開標本の機構，測定経路，出力・速度・構造変化を結び付けた追加評価を示す．結果は評価した標本・入力・条件に限られる．
\end{abstract}
\begin{eabstract}
We evaluate obfuscation in public WebAssembly mining artifacts by relating transformation scope to validated hashing throughput. After comparing prior evaluation targets and conditions, we check public test vectors and repeatedly measure ten distributed Proof-of-Theft artifacts, one SWAMPED input, and two locally generated transformations using public code. For xmr-wasm, three measurements of at least 30 seconds yield median rates of 9.11 hashes/s for the distributed reference and 6.18 hashes/s for the flattened artifact. Instruction growth and throughput loss rank differently for flattening and substitution. Some statically confirmed transformations occur on paths skipped by the measured v0 inputs, limiting explanations based on whole-file growth. Under the same assembly conditions, SWAMPED NOP insertion and integer constant splitting yield instruction-count ratios of 1.097 and 1.034 and median throughput ratios of 0.960 and 0.973, with overlapping observed ranges. Build-label discrepancies, optimization differences, and reused artifacts limit causal and implementation-independence claims. Additional measurements run locally without external pools. The contribution is an artifact-specific correspondence among mechanisms, measured paths, validated outputs, throughput, and structural changes, rather than a general description of obfuscation techniques. Findings are limited to the evaluated artifacts, inputs, and conditions.
\end{eabstract}
\maketitle

\section{はじめに}
WebAssembly（以下Wasm）は，ブラウザなどで動作するプログラムを表すバイナリ形式であり\cite{webassembly}，暗号資産の採掘にも用いられる．採掘はプログラムが行う処理を指す．一方，クリプトジャッキングは利用者の同意なく計算資源を採掘に使う行為であり，採掘機能があることとは区別する必要がある．採掘の調査や検出に関する研究では，収集したファイルに加え，評価のために加工した採掘プログラムも公開されている\cite{konoth2018,harnes2024,kim2025}．

難読化手法と検出への影響は既存研究で説明・評価されている\cite{harnes2024,cabrera2023,cao2023}．本稿で問うのは，公開された採掘標本で実際にどの処理が変換され，その処理を通る入力のハッシュ計算が正しいか，どのような性能変化が観測されるかである．設定名だけでは，適用範囲，通常処理との違い，実際の測定経路を確定できない．既存の手法説明を繰り返すのではなく，公開標本の根拠と出力・性能の結果を対応付けるため，次の二つを研究課題とする．
\begin{description}
\item[RQ1] 公開WebAssembly採掘データセットにはどのような難読化手法と実装上の仕組みが確認され，通常処理・研究用計測コードとどのように区別できるか．データセット・実装ごとにどのような傾向があるか．
\item[RQ2] 難読化はハッシュ出力の正しさ，計算速度，コード構造・命令数・ファイルサイズにどのような影響を与えるか．その傾向は複数の採掘実装で共通するか．
\end{description}
RQ1では複数の公開資料を調べ，どのような手法が含まれるかを整理する．その上で，代表的な採掘プログラムの命令列から仕組みを確認する．RQ2では，基準版と変換版を対応付け，静的な変化量に加えて，ハッシュ出力の照合と同一条件の速度測定を行う．RQ2の対象や条件を変更する場合は，その標本の通常版との対応，機構，適用範囲，計測処理をRQ1でも再確認し，両RQの結果と考察を更新する．一般用途のプログラムは，難読化についての判断が妥当かを確かめるために使い，採掘プログラムで手法を確認した例には数えない．

本稿の追加評価は，公開標本における変換の適用範囲と測定経路の対応，公表テストベクトルで出力を確認したハッシュ計算性能，静的な増加量と速度差の関係である．Proof-of-Theftでは配布設定の不一致を検出し，手法別の因果比較に使える範囲を限定した上で反復測定を行った．SWAMPEDでは空のimportを用いる既存評価と区別し，完全なハッシュ計算を復元した入力に公開変換コードを適用して同条件の対照を測定した．同じSHAの再収録や共通ソースの再利用も照合し，データセット名の数を独立した実装数と扱わない．新たな難読化アルゴリズムと検出回避能力は評価せず，未報告とする範囲は確認した先行資料に限定する．

\section{関連研究}
MineSweeperは，ブラウザ内での採掘の調査と検出を扱う\cite{konoth2018}．WasmBenchは実際に使われているWasmを集めたデータセットであり\cite{hilbig2021}，本研究では一般用途のプログラムの構造との比較に用いる．WasmGuardは，バイナリファイルを入力して悪性Wasmを検出する方法を扱う\cite{sun2025}．本稿で用いる悪性・良性の区分は配布元によるものであり，すべてのファイルに採掘機能があるかを本研究で確認した分類ではない．

Cryptic Bytesは，ソースコード，コンパイラ内部で使うLLVM中間表現，Wasmの各段階で変換を行い，難読化と採掘の検出との関係を評価した\cite{harnes2024}．本研究では，公開された通常版と変換版を対応付けて比較する．バイナリのサイズなど，同研究と重なる測定は公開ファイルの再測定に当たる．本研究では採掘プログラムを取り出し，手法ごと・元のプログラムごとの変化量を，ファイル内部の領域であるセクションの変化や，命令列で確認した仕組みと結び付けて説明する．サイズの差を測ること自体を新しい成果とはしない．

SWAMPEDは，変換後も検出器が検出性能を保てるかを体系的に評価している\cite{kim2025}．同論文で生成したデータセットと，本研究で入手した変換元のファイルやツールの使用例では，含まれるファイルの範囲が異なる．Proof-of-Theftは，実行中のデータの流れをグラフで表して採掘を識別する方法を扱い，難読化した採掘プログラムを配布している\cite{sermchaiwong2026}．本研究では，これらの公開ファイルの難読化機構を照合し，ハッシュ計算を最後まで実行した出力と速度を追加評価する．

WASMixerはバイナリ段階での難読化\cite{cao2023}，SEISMICは採掘を監視するための計測処理を扱う\cite{wang2018seismic}．両者の公開資料は，変換や計測処理によって命令列にどのような違いが生じるかを確認する補助資料とした．本文を確認したWASMixerのarXiv第1版の第5.2・5.5節では，ブラウザとJavaScriptへの依存を理由に採掘データセットを実行時間評価から除外している\cite{cao2023}．ESORICS 2024の刊行版も存在するが\cite{cao2024published}，評価本文の版間照合は未完了である．本稿のSWAMPED入力と同一SHAのファイルはWASMixerにも再収録されているが，今回の自生成変換はSWAMPEDの手法であり，WASMixerの変換手法を実行評価した結果とは区別する．JavaScriptの処理をWasmへ移すWobfuscator\cite{romano2022}のような言語をまたぐ変換は，本稿で調べるWasm内部の変換とは対象範囲が異なる．

\subsection{既存の性能評価と本研究の追加評価}
Cryptic BytesとWebAssembly diversification for malware evasionはハッシュ速度を評価しており，その対象を初めて性能評価したとは主張しない\cite{harnes2024,cabrera2023}．MineSweeperの表10はWasmとasm.js，CPU利用率による速度差を扱う\cite{konoth2018}．Proof-of-Theftの第6節はグラフの簡約と検出を評価し，確認した本文と補足資料には難読化前後のハッシュ速度と出力の比較結果がない\cite{sermchaiwong2026}．ただし公開コードにhashrate.jsがあるため，過去に測定自体が一切行われていないとは断定しない．SWAMPEDの第IV.D節は実行時間を評価するが，JavaScriptの処理がないためimportを空の関数で置き換える\cite{kim2025}．この測定と，正しいハッシュを最後まで計算した速度を区別する．MinerRayの第6.4節と表5の時間は検出用グラフの構築と走査に要する時間であり，採掘関数のhashes/sとは比較できない\cite{romano2020minerray}．本研究では，これらの公開標本の出力検証と実測を静的構造変化と結び付ける．別名のデータセットでも同一バイナリや同一ソースの再利用を調べる必要がある．


\begin{table*}[t]
\centering\small
\caption{確認した先行評価と本研究の位置付け．未確認は調査した本文・補足・公開資料で報告を特定できないことを指し，世界で未評価という意味ではない．}
\label{tab:prior-performance}
\begin{tabular}{p{0.22\textwidth}p{0.36\textwidth}p{0.34\textwidth}}
\hline
研究 & 確認した既存の評価 & 本研究で扱う範囲 \\
\hline
Cryptic Bytes\cite{harnes2024} & 難読化後のハッシュ速度（100秒） & 既存評価済みの比較対象．静的機構と構造変化を照合．\\
WebAssembly diversification\cite{cabrera2023} & 通常版6・変換版120，約1000秒とプールによる検証 & 評価設計の参照．現在の追加実測は公表入力とローカル評価．\\
MineSweeper\cite{konoth2018} & 表10のWasm/asm.jsとCPU利用率による採掘速度 & 評価済みの参照資料．配布物全件の難読化効果とは区別．\\
Proof-of-Theft\cite{sermchaiwong2026} & 第6節のグラフ・検出評価．出力と反復hashes/sの対比較報告は未確認 & 配布標本2実装の出力・速度・機構を追加照合．hashrate.jsの存在も記録．\\
SWAMPED\cite{kim2025} & 第IV.D節の空のimportを使った実行時間評価 & 必要なホスト処理を復元した完全なハッシュ計算と，自生成変換対を評価．\\
MinerRay\cite{romano2020minerray} & 第6.4節の検出用グラフ構築・走査の時間とメモリ評価 & 同一SHAの入力をSWAMPED側から実測．独立した新標本とは数えない．\\
WASMixer（arXiv第1版）\cite{cao2023} & 第5.2・5.5節で採掘D3を実行時間評価から除外 & 再収録入力への完全な計算を追加．今回の自生成対はSWAMPEDの手法．\\
\hline
\end{tabular}
\end{table*}

\section{対象と分析方法}
\subsection{対象の選定と配布物の役割}
元の論文や配布元の資料から採掘研究との関係が分かり，Wasmファイルを入手できた資料を主な対象とした（表\ref{tab:scope}）．Cryptic BytesとProof-of-Theftは，研究用に変換したファイルを調べる資料として扱う．MineSweeperとSWAMPEDのCryptojackersは，採掘に関連して収集されたファイルを調べる資料として扱う．WasmMalも関連するデータセットとして参照するが，悪性と分類されたファイルのすべてを，クリプトジャッキングに使われたと確認済みのファイルとはみなさない．
\begin{table*}[t]
\caption{分析対象の役割と確認範囲}
\label{tab:scope}
\centering\small
\begin{tabular}{p{29mm}p{51mm}p{78mm}}
\hline
資料 & 本稿の対象と規模 & 根拠と比較上の制約 \\
\hline
Cryptic Bytes & 採掘10プログラム，2配布条件，20,460件の比較 & 基準版を含む．16種類の設定とwasm-mutateによる変換．元プログラムと同条件の基準版を対応付ける．\\
Proof-of-Theft & 採掘プログラムの難読化用34ディレクトリ & 9種類の設定名に整理．元ファイル，計測版，中間ファイルを含むため，全件を難読化されたファイルと数えない．\\
MineSweeper & miner samples 34種類，解析成功33 & 配布元の採掘資料．加工元・適用手法は各ファイルで確定するわけではない．\\
SWAMPED & Cryptojackers 43種類，解析成功42 & 変換元として配布されたファイルを対象とする．他の悪性31種類とツール例9種類は補助資料．変換済みデータセット全体は未解析．\\
WasmMal & malware 7,512種類，解析成功7,510 & 配布元の悪性分類．各ファイルの採掘機能と難読化手法は全ファイルでは確認していない．\\
比較用資料 & WasmMal benign，WasmBench，SEISMIC，WASMixer & 一般的な構造と計測の確認に利用．採掘プログラムの定量比較には混入させない．\\
\hline
\end{tabular}
\end{table*}
Cryptic Bytesでは，配布先のパスと公開された区分から，CryptoNight系列の採掘プログラム10本を選んだ．配布条件は未最適化と最適化済みの二つである．各条件には，変換前の基準版10件，LLVM系で変換した80件，Tigressで変換した80件，wasm-mutateで変換した10,060件が含まれる．計20,460件を同じ条件の基準版と対応付けた．この中には，基準版をそれ自身と比べる20件も含む．また，元の10本には同じ系列から派生したプログラムが含まれ，独立した10種類のマルウェア系統を表すものではない．

Proof-of-Theftでは，ディレクトリ名に記された変換設定を整理した．ただし，フォルダにはビルド途中の\texttt{.o}ファイルや，計測処理を加えたWasmも含まれる．例えば，公開MakefileではWasabiを適用した後，入力ファイルを\texttt{.orig.wasm}に改名し，計測版をコピーする．このため，origという名前だけで難読化前の基準版とは判断できない．本稿ではこの資料をRQ1の設定一覧に用いる．ファイルの役割と条件をそろえられていない比較は，RQ2で測る難読化の変化量には含めない．

\subsection{実行せずに調べる特徴と難読化の仕組み}
入力ファイルを識別するためにSHA-256を保存し，WABT 1.0.42で検証と逆アセンブルを行い，セクションを解析した．ファイルは実行しない．命令数はコードに書かれた命令の個数であり，実行時に繰り返し実行される回数ではない．抽出する特徴は，関数数，複数の分岐先から一つを選ぶ\texttt{br\_table}，その分岐表がループ内にあるかどうか，直接・間接の関数呼出し，何もしない命令\texttt{nop}，メモリを読み書きするロード・ストア命令などである．ロード・ストアの命令数も，実行時のアクセス回数や読み書きするバイト数を表すものではない．また，名前の情報を持つnameセクションの有無と，命令を持つcode，データを持つdata，追加情報を持つcustomの各セクションの大きさを調べる．検証に失敗したファイルは，特徴がないとは判断せず，解析できなかったものとして分ける．

分析では，「配布元がどの変換を指定したか」「コードにどの特徴があるか」「命令列からどの仕組みを確認できたか」を区別する．平坦化では，次の処理を決める状態の値について，初期化，読出し，更新，共通の選択処理に戻る流れを追う．仮想化では，独自の命令列をどこから読み，読み出した命令に応じてどの処理を選ぶかを調べる．演算式の置換では入力値と計算結果の関係を，関数呼出しでは呼出し先を指定する番号の取得元を確認する．分岐表やXORがあるだけで難読化手法を確定することはしない．

通常処理と計測コードの対照には，旧予備解析の命令位置と変換前後の対応が残る例を再利用する．その結果と採掘標本への適用範囲は第\ref{sec:rq1-controls}節に示す．

難読化の仕組みを確認するため，元のCryptoNightを表すcn-0を選び，基準版と6設定の変換版を用意した．本稿では，根拠を確認した4設定を詳しく説明する．選んだファイルのうち，算術式の複雑化と定数の難読化については，配布元の設定とは別に命令列から仕組みを確認する作業が終わっていない．そのため，変化量を比較する際は配布元が示した設定名として扱う．詳しく調べたのは代表的な箇所に限られ，すべての変換版の全関数について処理の意味を解析したわけではない．

\subsection{変換前後の比べ方と結果のまとめ方}
命令数やファイルサイズを比べる際は，元のプログラムと最適化の条件をそろえる．元のプログラム$p$，配布条件$c$における基準版の値を$X_{p,c,0}$とする．手法$m$で作った変換版$v$の値を$X_{p,c,m,v}$とし，基準版からどれだけ増減したかを差分で，何倍になったかを比で求める．
\[
\Delta X=X_{p,c,m,v}-X_{p,c,0},\qquad
r_X=\frac{X_{p,c,m,v}}{X_{p,c,0}}.
\]
変化率は$100(r_X-1)$\%である．基準値が0の場合は比を定義できないため，差分を残す．集計は二段階で行う．まず，同じ元プログラムに同じ手法を適用して作ったファイルについて，変化量の中央値を求める．次に，元プログラム10本の値から，中央値，四分位，最小値，最大値を求める．大量に作られた変換版を，互いに独立した事例として統計検定に使うことはしない．示す分布は今回のデータの範囲を記述するもので，より広い集団についての信頼区間ではない．

LLVM系とTigressは，それぞれ8設定を分けて集計する．wasm-mutateは，個別の変換と変換を繰り返したファイルが混在するグループとして別に示す．このグループの中央値を，一つの手法の効果とは解釈しない．未最適化版と最適化済み版は，それぞれ同じ配布条件の基準版と比較する．本研究で新たに最適化を実行した実験ではない．元論文に記されたEmscripten 3.1.35，Tigress 3.3.2，LLVM 16系のemcc-obf，wasm-tools 1.0.33という環境に加え，入力パス，ハッシュ，取得元のリビジョンを記録した．すべてのバイナリをビルドし直して同一のファイルができることを確認したわけではない．

\section{RQ1：難読化手法とデータセットの傾向}
\subsection{配布物に含まれる手法の範囲}
表\ref{tab:settings}に，研究用に変換されたデータの設定を示す．テイント解析は，データがプログラム内でどのように伝わるかを追跡する解析である．算術式の複雑化は，計算結果を保ちながら式を複雑にする変換を指す．リテラルは，数値や文字列など，コードに直接書かれた値である．Cryptic Bytesでは，処理の流れ，算術式や定数，関数呼出しなどに対する幅広い変換を比較できる．個々の設定では，採掘プログラム10本を未最適化と最適化済みの二つの条件で比較できる．wasm-mutateには変換を繰り返して作った多数のファイルが含まれる．ファイルの収録数は，実際の攻撃でその手法が使われる頻度を表すものではない．
\begin{table*}[t]
\caption{研究用配布物の変換設定（すべてのファイルで仕組みを確認したものではない）}
\label{tab:settings}
\centering\small
\begin{tabular}{>{\raggedright\arraybackslash}p{32mm}p{94mm}p{32mm}}
\hline
資料／ツール & 設定名から整理した手法 & 本稿の収録範囲 \\
\hline
Cryptic Bytes／LLVM & 偽の制御フローの追加，平坦化，定数の難読化，ラッパー関数，間接分岐，基本ブロック分割，文字列暗号化，命令置換 & 各設定10本$\times$2条件 \\
Cryptic Bytes／Tigress & エイリアス解析への対策，テイント解析への対策，算術式の複雑化，リテラルの難読化，平坦化，ランダム関数生成，関数分割，仮想化 & 各設定10本$\times$2条件 \\
Cryptic Bytes／wasm-mutate & 複数種類の変換と繰り返し変換したファイル & 10,060件$\times$2条件 \\
Proof-of-Theft／emcc-obf & 偽の制御フローの追加，平坦化，分割，命令置換 & 各6設定ディレクトリ \\
Proof-of-Theft／Tigress & 算術式の複雑化，平坦化，分割，分割と平坦化の併用，仮想化 & 順に3，1，1，4，1ディレクトリ \\
\hline
\end{tabular}
\end{table*}
Proof-of-Theftでは，34ディレクトリを9設定に整理した．emcc-obfの4設定にはそれぞれ6つのプログラム名があるが，Tigressでは設定ごとの収録数が均等ではない．これは取得した配布物の構成であり，元論文の評価対象数や，完成した難読化ファイルの数と同じとは限らない．したがって手法間の効果を比較する際には，変換したプログラムや追加した計測処理の違いが，結果に含まれる可能性がある．

収集されたファイルでは，コードの特徴を抽出できても，どのような加工を受けたかが分からない場合がある．表\ref{tab:features}は，静的解析に成功したファイルのうち，各特徴が見られた割合を示す．採掘関連の2グループでは，解析に成功した全ファイルに間接呼出しがあった．ただし，これは全件で呼出しを間接化する難読化手法を確認したという意味ではない．比較用のグループにも，同じ構造が広く見られる．
\begin{table*}[t]
\caption{構造上の特徴の分布（分母は各群の解析成功数，単位\%）}
\label{tab:features}
\centering\small
\begin{tabular}{lrrrrr}
\hline
配布物の群 & 成功数／取得種類数 & ループ内分岐表 & 間接呼出し & nop & 手法の確認率 \\
\hline
MineSweeper／miner samples & 33／34 & 33.33 & 100.00 & 54.55 & 未推定 \\
SWAMPED／Cryptojackers & 42／43 & 42.86 & 100.00 & 57.14 & 未推定 \\
WasmMal／malware（参考） & 7,510／7,512 & 46.07 & 98.88 & 59.57 & 未推定 \\
WasmMal／benign（比較） & 7,510／7,512 & 52.77 & 81.85 & 20.27 & 未推定 \\
WasmBench／filtered（比較） & 8,449／8,460 & 56.66 & 82.73 & 25.79 & 未推定 \\
\hline
\end{tabular}
\end{table*}
各グループ内では，SHA-256が一致するファイルを一件にまとめた．SWAMPEDの取得済み採掘入力42種類のうち35種類は，保存済みの他データセットの取得記録にも同一SHAで含まれた．グループをまたいだ重複は残るため，表の各行を足して独立した利用例の総数とはしない．以前に探索用として選んだWasmMalとSWAMPEDの8ファイルでは，確認した分岐が入力や型に応じて変わり，その箇所を平坦化と確定できなかった．ファイル全体に難読化がないことを示す結果ではない．そこで以下では，変換前後を比較できる採掘プログラムを用いて，難読化の仕組みを説明する．

\subsection{通常処理と研究用計測による対照}
\label{sec:rq1-controls}
WasmMalでは，malwareに分類された121件，benignに分類された3件に，SEISMICの計測用関数と一致する関数があった．一般用途のA-StarとCreaturePackでは，カウンタ更新の命令列を取り除くと，元の全147関数，全756関数と命令列が一致した．命令数の増加だけでは難読化と計測を区別できない具体例であり，RQ2ではWasabiのhookを持つ標本を混ぜずに測定する根拠になる．ただしWasabiとSEISMICが同じ計測機構であるとは扱わない．

CreaturePackの関数400では，引数が指す位置から読んだ値と74のANDを分岐表の選択値に用い，各分岐は8，16，0，10を返す．関数には状態を書き換えるストアもループもなく，シンボル名も数値の基数選択を支持する．これに対してWASMixerのbase64-cliの関数515では，元の3ストアを含む命令列を保持しながら，状態0でその処理へ分岐し，状態1へ更新してループへ戻り，次の選択で終了する．この対照は分岐表の存在より状態の生成・更新・戻り先を確認する必要を示す．

別のseqtk対照でも，状態0からcalloc(1,28)を呼び，結果を保存して状態1へ更新し，選択処理へ戻ってから結果を返す経路を確認した．元の関数11と45が同じ処理を持つため，唯一の対応関数は確定しない．これらは一般用途の選択例による局所的な対照であり，採掘標本の確認件数や，全件分類の精度を表さない．

\subsection{状態変数を用いる制御フロー平坦化}
平坦化では，状態変数に保存した値で次の処理を選ぶ．cn-0をTigressのflatten設定で変換した版では，外部から呼び出すための名前であるexport名\texttt{hash\_cn}に対応する関数33に，状態変数があった．関数内でメモリを使う際の基準位置に32を加えたアドレスへ，初期値5を保存する．ループ内でその値を読み出し，デフォルトの分岐先を含む12要素の\texttt{br\_table}で処理を選ぶ．各処理は状態を書き換えて，共通のループに戻る．同じ状態の保存先への書込みは，関数内で9箇所確認した．通常版で同じexport名に対応するのは関数53である．このように，関数番号が同じかどうかではなく，export名を使って対応を調べた．

この構造では，次に実行する処理はコードの並びだけでは分からず，保存された状態の値にも依存する．分岐表があるかどうかに加え，状態を作り，読み出し，書き換え，ループに戻る流れを追うことで，flatten設定が具体的にどのようなコードになったかを説明できる．ただし，すべての実行経路で元と同じ動作をすることを証明したわけではない．

\subsection{仮想命令列の読出しと処理の選択}
コードの仮想化（以下，仮想化）では，独自の命令列を読み出し，その命令に対応する処理を選んで実行する．Tigressのvirtualize設定で変換した版では，同じexport名に対応するのは関数56である．メモリを使う際の基準位置に24を加えた領域に，命令の読出し位置を指すポインタが保存され，その初期値は168944であった．ループ内では，ポインタが指す先から\texttt{i32.load8\_u}で1バイトを読む．そこから16を引いた値を使い，デフォルトの分岐先を含む233要素の分岐表で処理を選ぶ．処理の中には，ポインタを1または4進めて保存し，共通の読出しループに戻る箇所がある．

これは，命令列を読み出して対応する処理を行うインタプリタの構造と解釈できる．前節の平坦化が状態の値で処理を選ぶのに対し，この例では独自の仮想命令列と，そのどこを読んでいるかを示す位置の情報を持つ．ただし，すべての仮想命令の意味や，元の各命令との対応までは復元していない．RQ2では，この仕組みとcode・data両セクションの増加を比べる．

\subsection{算術式の置換と間接呼出し}
算術式の置換では，計算結果を保ちながら式を複雑にする．LLVMのsubobf設定で変換した版の関数23では，通常版の左シフト$x\ll3$が，次の32ビットの演算式に置き換わっていた．
\[
(x\mathbin{\&}\mathord{\sim}8)(\mathord{\sim}x\mathbin{\&}8)
+(x\mathbin{|}8)(x\mathbin{\&}8)
\]
ここで，隣り合う括弧は乗算を表す．式中の記号は，ビット単位のAND，OR，NOTである．$u=x\mathbin{\&}\mathord{\sim}8$，$v=x\mathbin{\&}8$とおくと，$x=u+v$，$v\in\{0,8\}$であり，式は$u(8-v)+(u+8)v=8x$となる．32ビットの範囲を超えた桁を捨てて計算すると，元のシフトと一致する．両版では，この値を対応する引数として関数24に渡す．確認したのはこの箇所で計算結果が同じになることであり，モジュール全体の動作が同じことではない．

エイリアス解析は，異なる参照が同じメモリ領域を指すかを調べる解析である．これへの対策を指定したTigressのantialiasanalysis版では，読み込んだ二つの値の比較結果から1436または8を選び，66688を加えたアドレスからテーブル内の番号を読み出して，間接呼出し命令\texttt{call\_indirect}に渡す箇所があった．この箇所では，呼出し先がメモリから得た番号に依存することを確認できる．ここで確認したのは変換設定版の呼出し先の取得方法であり，通常版の対応箇所との差分や呼出し先の同値性は未確定である．従って，この観測だけから難読化によって呼出し間接化が追加されたとは断定しない．間接呼出し命令の数と，呼出し先の番号をどこから得ているかは分けて記録した．

\subsection{実行評価対象での機構と配布設定の照合}
\label{sec:proof-mechanisms}
RQ2の実測対象について，配布設定と実際の命令列を照合した．Proof-of-Theftのwebminerpoolhashcnでは，配布基準と4変換ディレクトリのMakefileがすべて\texttt{-enable-bcfobf}を指定していた．設定と名称が一致せず，過去の生成指示も確定できないため，後節の速度は配布標本間の比較とし，4手法のオーバーヘッドには用いない．xmr-wasmでは名称に対応するフラグがあり，5標本に同梱されたC/C++・ヘッダ46ファイルは相対パスとSHA-256で基準版と一致した．配布ソース上の対応は確認できるが，過去のコンパイラ，入力，オブジェクトファイルまで同一だったとは保証しない．

\paragraph{平坦化の機構と適用箇所}xmr-wasmの公開hash関数では，通常版のnonce・target処理に対し，状態を読み出して次の処理を選び，状態を更新して共通ループへ戻る制御を確認した．この関数には\texttt{br\_table}がなく，分岐表の存在を必要条件としない具体例である．8・16バイトのXOR補助関数112・118でも，条件検査，XORと書込み，添字更新，終了を状態で表すように変わった．通常版との引数，上限8・16，増分1の対応を確認した．コンテキスト初期化83，鍵取り込み82，解放86にも同じ機構があり，それぞれ4・24・8状態をレビューした．引数検査や確保・コピー・解放と，同梱OpenAESソースの処理が対応するが，全入力の同値性や過去のビルドとの同一性の証明ではない．

\paragraph{呼出し位置と今回の測定経路}通常版の関数111では，補助関数118の2呼出しが主ループ内にある．一方，初期化・鍵取り込み・解放の呼出しは構造化ループ外にあり，ハッシュごとの処理と主ループ内の処理を区別する．関数112の最初の呼出しはvariant 1の条件内である．関数113でも平坦化の入口と最初の条件を確認し，初期化命令列239行を通常版と対応付けたが，全状態の意味は未確認である．同梱CryptoNightR用演算生成処理と特徴が対応し，呼出しにはvariantが4以上という条件がある．通常版と平坦化版の呼出しまでの命令列は局所変数名を対応付けると一致する．今回の3入力は先頭バイトが7未満でvariant 0を選ぶため，この呼出し経路には入らない．機構の存在と実測中の経路を分け，動的な回数や時間寄与を存在だけから決めない．各状態値，命令位置，対応するソース抜粋は保存した解析記録に残した．さらに関数111では，平坦化版の本体命令列の増加はvariant 4以上のブロック内に収まった．そのブロックを除く3,609行は，局所変数とラベルの名前を対応付けると命令，数値オペランド，参照する局所変数の型が一致した．これは本体内の比較であり，呼出し先の変化やモジュール全体の同値性を否定・証明するものではない．

命令置換版の公開hash関数にも，元の式との対応を確認した．入力先頭のバイト$x$からvariantを求める式$x-6$は，定数$C=1368051000$を用いた$(x-C)-6+C$に置き換わる．nonceの更新$n+1$は$n-(0-1)$に置き換わる．どちらもWasmの32ビット整数の剰余演算で元の式と一致し，単に演算命令の多さから置換と推定した例ではない．ただし今回のv0入力では先頭バイトが7未満なので前者の経路を通らず，最大targetでnonceの変化がないため後者の経路も通らない．この機構確認と，後節の測定中に実行された処理を区別する．さらに公開hashから関数35を介して呼ぶ関数111でも置換を確認した．通常版ではオフセット1048のループ添字$x$を4ビット左シフトしてアドレス計算に用いるが，置換版では$(x\mathbin{\&}\mathord{\sim}16)(\mathord{\sim}x\mathbin{\&}16)+(x\mathbin{|}16)(x\mathbin{\&}16)$を用いる．$u=x\mathbin{\&}\mathord{\sim}16$，$v=x\mathbin{\&}16$とおくと，$v\in\{0,16\}$なので$u(16-v)+(u+16)v=16x$となり，32ビット剰余演算で元のシフトと一致する．この箇所は上限16384と8を持つ二重ループ内にあり，ラッパ以外にも機構の根拠がある．呼出し経路，命令位置，式の導出と2061入力の局所演算確認を保存した．動的な経路被覆とモジュール全体の意味保存は未確認である．

分割設定版では，関数35のスタック確保が通常版の32バイト一括から16バイトずつ3回の計48バイトへ変わり，5引数のうち2引数の保存位置も変わった．保存済み命令列の局所的なデータフローを追うと，両版は関数111へ同じ6引数を渡し，復帰時には入口のスタック位置を復元する．しかし呼出し時のスタック位置と書込み先は異なり，全体の同値性は証明していない．このラッパーには追加の分岐や関数呼出しがない．補助関数112・118でも一時保存先が変わり，外側のblock内にloopを置く通常版から，loop内のblockを抜けて終了する表現へ変わる．XORの対象，上限8・16，増分1とループへの戻りは対応する．これらの関数には新たな呼出しやdispatcherはなく，WasmのblockはLLVMの基本ブロックとも異なるため，ここで確認した変更だけを基本ブロック分割や関数分割の根拠とはしない．

emcc-obfの公開説明では，\texttt{enable-splitobf}はLLVMの基本ブロック分割を指定する\cite{emccobfdocs}．関数分割やスタック領域の分割とは区別する．配布Makefileの\texttt{split\_num=3}と公開説明は対応するが，当時の変換前後のLLVM中間表現がないため，Wasmで見える変更から適用履歴を確定していない．

通常処理の対照として，xmr-wasmの配布基準にも命令列を解釈する構造がある．関数111では，opcode，宛先・参照元の番号を読み出す処理と70箇所の分岐表が，同梱CryptoNightRの\texttt{v4\_random\_math}による演算選択の展開と対応した．最初の選択と乗算処理を通常版・平坦化版で確認すると，分岐表が比較と条件分岐へ変わっていた．しかし状態を更新するdispatcherの追加とは区別する．暗号アルゴリズム自身が命令列を解釈する場合もあり，この構造の存在だけを難読化の仮想化の根拠にはしない．全70箇所の各処理の同値性は未確認である．cryptonight-hashの配布基準にも，内部状態の下位2ビットで4種類の最終ハッシュを選ぶ間接呼出しがある．同梱ソースのBlake，Groestl，JH，Skeinによる最終処理と，内部関数の200バイトの状態および出力ポインタを渡す命令列が対応するため，間接呼出しの存在だけを難読化による呼出し変更と判定しない．通常版と5変換版では，初期データの表番号1～4が同じ4関数を指し，通常版の引数順，バイト長からビット長への変換，256ビットの指定が同梱ソースの4処理と対応した．ただし，計算中に表が変わらないことや，各暗号処理の本体とモジュール全体の同値性は未検証である．

この比較は，通常の数値処理に分岐表が現れる例と，平坦化でも分岐表が現れない例を合わせて，静的な指標の有無だけでは判断できないことを示す．またWasabiのhook importを持つ版を除外し，非計測版の機構を確認することで，計測の追加処理を難読化の根拠に混ぜない．旧予備実験のA-StarとCreaturePackの命令列復元は，計測コードとの区別を確認する対照として再利用する．

追加候補のcryptonight-hashでは，内部ハッシュ候補関数10の平坦化を静的に確認した．状態値による条件選択，計算本体，カウンタ更新の3状態が共通の選択ループへ戻り，通常版の上限524288を持つ反復処理の一部に対応する．両版のカウンタ初期化は0，継続条件は符号なしの上限比較，更新は1加算である．通常版のループへの分岐に対し，平坦化版は次状態を保存して選択ループへ戻る．本体の157行は，局所変数名とスタック保存位置の16バイトの差を対応付けると命令列が一致し，6箇所の直接呼出しも同じ順序である．ただし，呼出し先によるメモリ更新を含むループ全体の同値性と実行経路は未確認であり，この候補のハッシュ出力・速度は未測定である．後節の13バイナリの実測結果には含めない．

\paragraph{RQ2の標本との対応}追加のローカル反復測定の13バイナリ（Proof-of-Theft配布10，SWAMPED入力1，自生成変換2）をSHA-256で照合し，配布設定，確認した機構，適用範囲，計測hookの有無，出力照合と速度の根拠を同じ一覧へまとめた．SWAMPEDの自生成対はnop挿入30関数3089箇所，整数定数分割29関数542箇所で，挿入命令の除去と定数式の復元により命令列全体が基準版に戻ることを確認した．これは元配布の変換版の機構を確認した結果ではない．一方，webminerpoolhashcnの設定不一致，xmr-wasmの分割・偽制御フローの機構未確認は一覧にも残し，出力一致だけで難読化ラベルを確定しない．webminerpoolhashcnの公表v0ベクトル照合と，速度測定で用いたcn-picoの受理済み出力照合はアルゴリズムを分けて記録する．

\subsection{RQ1への回答}
研究用の公開ファイルには，処理の流れの平坦化や分割，算術式の複雑化や置換，関数呼出しの変更，仮想化などの変換設定がある．変換設定の種類や適用対象は資料によって異なる．代表的な採掘プログラムでは，状態の値による処理の選択，仮想命令の読出し，計算結果を保つ算術式の置換を確認した．メモリから得た番号による間接呼出しも観測したが，通常版との対応が未確定の箇所は，難読化で追加された機構とは扱わない．cryptonight-hashの通常版にも最終ハッシュの選択による間接呼出しがあり，その存在だけでは難読化と区別できなかった．RQ2の追加対象では，xmr-wasmの公開hashと一部の呼出し先で平坦化・置換を確認し，SWAMPED入力への自生成変換ではnop挿入と整数定数分割の変更位置を確認した．一方，配布設定の不一致や分割・偽制御フローの機構未確認を残した標本は，確認済み手法の性能比較とは区別する．収集されたファイルにも関連する構造は見られたが，各手法を確認できた割合は推定していない．したがって，ここで説明できるのは公開資料の設定と確認した仕組みの傾向であり，攻撃全体で各手法が使われる割合ではない．

\section{難読化に伴うコード構造・サイズの変化}
\subsection{手法別の命令数とファイルサイズ}
RQ1で確認した仕組みを踏まえ，ここでは難読化によってコードの量と構造がどれだけ変わったかを調べる．図\ref{fig:methods}に，採掘プログラム10本の変化を示す．20,460件の比較では，すべてのファイルから特徴を取得できた．基準版をそれ自身と比べた20件では，差分が0となった．また，サイズと命令数の比が保存済みの対応表と一致することを確認した．図の手法名は配布元の設定名であり，全ファイルの仕組みを独自に調べて分類したものではない．
\begin{figure*}[t]
\centering
\includegraphics[width=\textwidth]{../results/mining-methods/figures/method-size-instructions.pdf}
\caption{手法別の命令数比とファイルサイズ比．各配布条件の基準版との比較．元プログラム10本の値をまとめた統計であり，独立した10系統のマルウェアについて推定したものではない．手法名は配布元の設定．}
\label{fig:methods}
\end{figure*}
未最適化の配布版では，Tigressの算術式の複雑化設定encodearithmeticによる命令数比は，中央値が2.515，最小値から最大値までの範囲が2.202--2.773だった．サイズ比は2.601で，個別の設定の中でも大きかった．LLVMの命令置換設定subobfでも，命令数比は1.461，サイズ比は1.477となった．subobfでは，前節で示したように一つの演算が複数の論理・算術命令に置き換わる．この例は命令数が増える仕組みを具体的に示すが，増加したすべての命令をこの一例で説明するものではない．

同じ平坦化でも，未最適化版の命令数比は，LLVMのcffobfが1.098，Tigressのflattenが1.394だった．最小値から最大値までの範囲は，それぞれ1.088--1.439，1.387--1.407であり，中央値とばらつきが異なる．したがって，手法名が同じなら命令数も一定の倍率で増えるとはいえない．この比較には，変換する段階や実装方法，適用範囲の違いも含まれるため，一方が難読化手法として優れているとは判断しない．

最適化済み版でも，encodearithmeticの命令数比は1.655，subobfは1.327となり，同じ条件の基準版との差が残った．一方，LLVMのfuncwraは，未最適化で1.006，最適化済みで1.005と増加が小さい．これらは配布版どうしを比較した結果であり，追加されたすべての命令が最適化で除去できないことを意味しない．複数の変換が混在するwasm-mutateでは，未最適化と最適化済みの順に，命令数比の中央値が1.395と1.350，サイズ比の中央値が1.396と1.330だった．ただし，個々の変換がどれだけ増加に影響したかは分けて調べていない．

\subsection{制御フロー，呼出し，メモリ操作}
命令数の変化の内訳を調べるため，表\ref{tab:structure}に代表的な設定での構造の差を示す．平坦化では，未最適化・最適化済みの両条件で，ループ内分岐表が中央値で70増えた．状態をメモリに保存する実装に対応して，書込みを行うストア命令も増えた．エイリアス解析への対策では，未最適化版で直接呼出しが218減り，間接呼出しが358増えた．減少数と増加数は一致せず，関数の生成などもあるため，直接呼出しが間接呼出しに一対一で置き換わったとは結論しない．
\begin{table*}[t]
\caption{代表設定の構造変化（同条件の基準版との差，10プログラムの中央値）}
\label{tab:structure}
\centering\small
\setlength{\tabcolsep}{4pt}
\begin{tabular}{llrrrrrr}
\hline
配布条件 & Tigress設定 & 関数 & ループ内分岐表 & 直接呼出し & 間接呼出し & ロード & ストア \\
\hline
未最適化 & エイリアス解析への対策 & 33 & 0 & $-218$ & 358 & 2,361 & 6,290 \\
 & 算術式の複雑化 & $-1$ & 0 & $-1$ & 0 & 12,834.5 & 5,997 \\
 & 平坦化 & $-1$ & 70 & $-1$ & 0 & 72 & 6,733 \\
 & 仮想化 & $-1$ & 91.5 & $-1$ & 0 & 1,150.5 & 6,967.5 \\
最適化済み & エイリアス解析への対策 & 69 & 0 & $-185$ & 358 & 2,201 & 6,297 \\
 & 算術式の複雑化 & 1 & 0 & 0 & 0 & 3,392 & 5,997 \\
 & 平坦化 & 1 & 70 & 1 & 0 & 101 & 6,743 \\
 & 仮想化 & 1 & 98 & 0 & 0 & 218.5 & 6,968.5 \\
\hline
\end{tabular}
\end{table*}
関数の数は，難読化によって必ず増えるわけではない．未最適化の平坦化・仮想化では中央値で1減り，最適化済み版では1増えた．これは，関数の処理を呼出し元に展開するインライン化や，補助処理などを含む全体の差である．差分のすべてを，新しく追加した補助関数の数とはみなせない．仕組みの解析では，export名と命令列を使って代表的な関数を対応付けた．すべての補助関数がどのように作られたかを特定したわけではない．

読込みを行うロード命令と，書込みを行うストア命令も，同じように増減するわけではない．平坦化では状態の保存に関係するストアが増えた．仮想化では，仮想命令の読出しや，仮想的なスタックの操作も加わる可能性がある．ここでの説明は代表的な命令列で確認した範囲に限られる．コードに書かれた命令の数から，実行時のアクセス回数やメモリ使用量は推定しない．

\subsection{仮想化のサイズ増加とセクション}
仮想化を行うvirtualizeの命令数比の中央値は，未最適化で1.417，最適化済みで1.418と近い．一方，ファイルサイズ比は1.834と2.322で異なる．ファイルには実行命令以外の情報も含まれるため，サイズの変化は命令数比だけでは説明できない．以下ではcn-0を例に，ファイル内のどの領域が増えたかを調べる．

cn-0の未最適化版では，基準版の310,503バイトに対し，仮想化版は573,954バイトで，263,451バイト増えた．このうち，命令を持つcodeセクションは165,827バイト，データを持つdataセクションは97,625バイト増えた．最適化済み版では，140,103バイトから324,681バイトに変わり，184,578バイト増えた．codeセクションの増加は86,969バイト，dataセクションの増加は97,598バイトだった．ファイルには他のセクションや長さを記録する情報もあるため，codeとdataの差だけを足してもファイル全体の差にはならない．

最適化済み版では増えたバイト数は少ないが，比較の基準となるファイルも小さいため，増加の倍率は大きくなる．dataセクションの増加量は両条件で近く，仮想命令列を持つ仕組みとも矛盾しない．ただし，増えたdataの全バイトを仮想命令として復元したわけではない．また，異なる指標の中央値を足しても，厳密なサイズの内訳にはならない．このため，ここでは一つの変換前後の組について測定した差を用いた．

\subsection{名前の情報と集計値の限界}
20,460件の比較では，基準版にも変換版にもnameセクションがなかった．そのため，名前を消す変換によってどれだけ変化したかは評価できない．ただし，nameセクションの情報と，外部から関数を呼ぶためのexport名は別である．export名までなかったわけではなく，仕組みの解析には\texttt{hash\_cn}を使った．nameセクションがないことを，意図的に名前を難読化した証拠とはしない．

中央値などの集計値が近くても，同じ変換が行われたとは限らない．Tigressのantitaintanalysisとrandomfunsでは，未最適化版の命令数比の中央値はともに1.332だった．しかし，同じプログラムを同じ条件で比較した20組のファイルでは，ハッシュが一致しなかった．命令数とサイズは変化の大きさを示すが，それだけで変換の仕組みを特定することはできない．

\subsection{静的な変化のまとめ}
手法によって，命令数とファイルサイズの増加量や，分岐表，関数呼出し，メモリを読み書きする命令の変化が異なった．算術式を展開する変換では命令数が大きく増え，平坦化では状態の値による処理の選択が加わった．エイリアス解析への対策では間接呼出しが増えた．仮想化ではdataセクションも増えるため，命令数比とファイルサイズ比は一致しない．傾向は配布条件と元のプログラムによって変わるため，どのプログラムでも同じ倍率で増えるとはいえず，実行時の性能を評価した結果ともみなせない．

\section{ハッシュ計算の正しさと性能}
\label{sec:runtime}
\subsection{環境と正しさの確認}
Windows上のLinuxコンテナで，Wasmを変更せず，公開ソースとJavaScriptに基づくホスト処理を用いた．Cryptic Bytesのcn-pico-trtl通常版とLLVMの偽制御フロー設定bcfobfを2 workerで実行し，本人の機器と公開受取アドレスによる通常採掘としてFastpoolへ接続した．xmrig-proxyは6.18.0である．元論文のCoinHive実装と同じランタイムではなく，1ハッシュごとのworker通信の負荷も含む．1,001サンプルから先頭10件を除き，991件と990差分区間を保存した．通常版で13件，変換版で10件のshareを提出し，すべての受理を記録した．各条件1回の試行であり，性能差の確定には反復が必要である．旧試行の読込み失敗やジョブ更新の不具合は，修正後の完了試行と区別して記録した．

Proof-of-Theftのwebminerpoolhashcnでは，配布されたcn.jsからminifyされたimport/exportと時刻関数の対応を確認した．元のhash\_cnで対応するcn-picoの引数を用いた．Wasabi計測版を除き，配布基準と，平坦化，分割，命令置換，偽制御フローというディレクトリ名の4標本を評価した．第\ref{sec:proof-mechanisms}節の不一致があるため，名前は変換の確定ラベルではない．別の採掘試行で受理された13件のジョブとnonceから入力を復元し，記録された出力を期待値とした．ネットワーク無効のコンテナでNode.jsから計算した結果，5標本すべてで13件全件が一致した．Proof-of-Theft側の標本を新たにプールへ提出した結果ではなく，記録済みの受理出力との照合である．有限の入力の結果であり，全入力の正しさは保証しない．

分割と平坦化を併用したTigress版では，hash\_cnの呼出しで間接呼出しに関する実行時エラーが生じた．ホスト実装，ビルド条件，標本のどこに原因があるかは未確定であり，難読化による機能破壊とは結論しない．仮想化版も，配布元のJavaScriptを用いた対照呼出しで同種のエラーとなった．アダプタだけに固有の不具合ではないことを確認したが，配布物のビルド・リンクの問題を含む原因の確定は未完了である．両者は照合失敗の理由とともに記録する．cryptonight-hashの7標本は，非計測版に入力ごとのハッシュ関数のexportがなく，配布JavaScriptにもその呼出し口を確認できなかった．付属ビルドスクリプトが指定するソースは，内部で84バイトの入力とnonceを生成して計算を繰り返すため，mainを呼ぶだけでは公表入力に対する出力照合にならない．6標本ではmainから呼ばれる内部関数に，入力・出力ポインタと長さに対応する可能性のある3引数を確認したが，関数の対応と出力検証は未完了である．別途，呼出し口を追加した評価用派生物を作る場合は，配布物そのものと区別する必要がある．この制約を出力不一致や，ハッシュ計算コードの不在と混同しない．

\subsection{公表テストベクトルによる別実装の照合}
独立した期待値としてXMRig v6.18.0の\texttt{CryptoNight\_test.h}に公表されたCryptoNight v0の入力5件と出力を用いた\cite{xmrigvectors}．webminerpoolhashcnの5標本は配布元のcn.jsから呼び出して5件全件が一致した．xmr-wasmは元ソースで入力先頭の値からvariantを選択するため，公表入力のうちv0に対応する3件を使用し，通常版と4変換版で全件一致した．後者の公開hash関数はtargetを満たすまでnonceを進めるが，最大targetを与え，確認した出力ではnonceが変化せず1回のハッシュに対応することも確かめた．デモのmainは実行せず，ネットワークも無効にした．これは有限のテストによる確認であり，全入力に対する意味保存の証明とは区別する．

\subsection{同一条件の短時間反復測定}
出力を照合できた5標本についてNode.jsの1 workerで同期的に計算し，1秒のウォームアップ後に5秒間を3回測定した．13入力を順に繰り返し，標本順を反復ごとに変えた．速度はハッシュ数を実際の経過時間で割って求め，中央値と最小・最大を示す．イメージID，Node.jsとV8の版，CPU型番，対象と測定コードのSHA-256，生データを保存した．
\begin{table}[t]
\caption{webminerpoolhashcn配布標本の予備測定（3反復，H/s．変換名はディレクトリ名で，適用手法の確定ではない）}
\label{tab:proof-preliminary}
\centering\small
\begin{tabular}{lrrr}
\hline
配布名 & 中央値 & 最小 & 最大 \\
\hline
配布基準 & 204.04 & 203.01 & 211.15 \\
平坦化 & 213.29 & 206.45 & 221.98 \\
分割 & 222.99 & 220.84 & 226.14 \\
命令置換 & 225.20 & 221.39 & 229.37 \\
偽制御フロー & 212.24 & 207.78 & 217.47 \\
\hline
\end{tabular}
\end{table}

この測定はブラウザと外部プールを使う長時間実験とは別条件である．配布名と設定の不一致に加え，ビルド・最適化条件と他のホスト処理の負荷を完全には分離していないため，速度比を難読化だけの因果効果とは扱わない．元標本同士の比較を予備結果として示し，時間を延ばした反復とビルド条件の対照を今後行う．最終出力の一致も各測定で確認した．


\subsection{xmr-wasmの反復測定と機構との対応}
\label{sec:xmr-runtime}
別の配布実装xmr-wasmの5標本を，配布元のJavaScriptから同期的に呼び出した．Node.js v18.20.4，V8 10.2.154.26-node.37，Intel Core i9-11900K上の隔離コンテナを用い，ネットワークを無効にした．各標本で公表入力3件の一致を確認した後，3秒のウォームアップと30秒以上の計算を3反復した．入力を順に循環させ，各反復で標本の実行順をずらした．スレッド数は1とし，実際の経過時間と完了したハッシュ数から速度を求めた．測定終了時の出力とnonceの不変も15ブロックすべてで確認した．表\ref{tab:xmr-runtime}に実測速度と命令数比を示す．3反復の範囲は信頼区間ではなく，ホストの他の負荷を完全には制御していない．
\begin{table}[t]
\centering\small
\caption{xmr-wasm配布標本の実測（各30秒以上，3反復）．命令数比は配布基準に対する値．偽制御フロー版は最適化条件が異なる．}
\label{tab:xmr-runtime}
\setlength{\tabcolsep}{3pt}
\begin{tabular}{lrrr}
\hline
配布名 & H/s中央値 & 最小--最大 & 命令数比 \\
\hline
配布基準 & 9.11 & 8.50--9.34 & 1.000 \\
平坦化 & 6.18 & 5.86--6.22 & 1.166 \\
分割（設定名） & 8.59 & 8.46--9.01 & 1.033 \\
命令置換（設定名） & 6.72 & 6.61--6.82 & 1.655 \\
偽制御フロー（$-O2$） & 23.08 & 22.69--24.05 & 0.413 \\
\hline
\end{tabular}
\end{table}

状態による処理選択を確認した平坦化版は，配布基準と比べて命令数1.166倍，サイズ1.170倍，速度中央値比0.678であった．これに対し，命令置換の設定を持つ標本は命令数1.655倍，サイズ1.641倍だが，速度比は0.738であった．従ってこの配布標本比較では，静的な増加量の順序と速度低下の順序は一致しない．置換版のhashラッパでは2つの同値式を確認したが，両式とも今回の測定では通らない経路にある．呼出し先の関数111のループ内にも同値なシフト式の置換を確認したが，変換範囲の全体と動的経路被覆は未確認であり，増加した全命令が実行されたとは扱わない．平坦化の機構はhashラッパ，XOR補助関数112・118，鍵取り込みの関数82，初期化の関数83，解放の関数86で確認したが，各関数の動的な呼出し回数は測定していない．その他の処理も測定時間に含まれるため，これらの機構だけから低下の全量を説明しない．関数113では平坦化の入口を確認したが，variant 4以上に限定された呼出しは今回のv0入力では通らず，その直接の実行負荷を速度低下の根拠にはしない．未実行コードがコンパイルや配置へ与える間接的影響は別途未検証である．通常版にある70箇所の分岐表もvariant 4以上の同じ条件内にあり，今回のv0ではその経路を通らない．従って分岐表の消失だけから，v0の速度低下を説明しない．関数111の本体命令列は，この未実行ブロック内で9,428行増え，それ以外は名前を対応付けると一致した．したがって，本体の増加分をそのままv0で追加実行された処理とは扱えない．補助関数の変換に加え，増えた局所変数の宣言やコンパイル・配置の違いを分ける追加評価が必要である．ここでの行数は保存WATの命令列の差であり，動的命令数やファイルのバイト増加量ではない．

各反復の変換版速度を同じ反復内の基準版速度で割ると，平坦化版は0.661--0.689，置換版は0.708--0.790で，3回とも基準を下回った．一方，分割設定版は0.906--1.011で差の方向が入れ替わった．これらは逐次測定した観測比の範囲であり，同時測定した対や信頼区間ではない．分割設定版の速度中央値比は0.943で，命令数比は1.033であった．第\ref{sec:proof-mechanisms}節ではラッパーとXOR補助関数の保存先・制御表現の変更を確認したが，分割の機構と全適用範囲は未確認である．この差を分割手法によるオーバーヘッドと確定しない．

偽制御フロー版は23.08 H/sで基準の9.11 H/sを上回り，命令数も0.413倍に減少した．しかし同梱Makefileでこの版だけに明示的な\texttt{-O2}があり，コンパイラの最適化による影響を含む．難読化が計算を高速化したとは結論しない．同じ最適化で再生成した対照がない現段階では，値は配布物の観測として示す．webminerpoolhashcnとの絶対速度の差についても，実装，ホスト呼出し，ビルド条件，測定時間が異なるため，難読化手法の優劣には用いない．

\subsection{SWAMPEDの公開入力標本での完全なハッシュ計算}
\label{sec:swamped-hash}
SWAMPEDのCryptojackers\_43にある取得済み42標本を調べ，ハッシュ関数のexportとホスト依存性を照合した．うちSHA-256が\texttt{046dc081}で始まる標本について，\texttt{\_cryptonight\_create}と\texttt{\_cryptonight\_hash\_variant\_0}を呼び出した．元のJavaScriptが付属しないため，Wasmのimportと命令列に基づいて旧Emscriptenのメモリ・表・時刻・UTC変換・プロセスIDを提供した．必要な処理を空の関数へ置き換えず，未対応のホスト呼出しはエラーにする．初期メモリは16 MiBで，ホスト側の配置は元JavaScriptと同一であるとは主張しない．計算後の出力はXMRigの公表v0入力5件すべてと一致した．

同じ隔離ランタイムで，1 worker，3秒のウォームアップ，30秒以上の3反復を行った．中央値は27.78 H/s，最小26.54 H/s，最大28.13 H/sで，全反復の最終ハッシュ出力も一致した．この標本は公開の変換入力であり，この測定単独では難読化のオーバーヘッドを示さず，後節で同条件の自生成変換対を別に比較する．既報のimportを空にした実行時間とは，ハッシュを最後まで計算して期待値を確認した点が異なる．ただし既報の測定条件とは一致せず，両者の数値を直接比べない．取得記録全体でSHA-256を照合すると，この入力はMinerRay，WasmBench，WASMixer，WasmMalにも同じバイナリとして収録されていた．初期の採掘候補一覧だけを用いた照合ではこの重複を見落としていたため，全取得記録との照合へ改めた．同じ標本の再収録を独立した実装とは数えず，本稿の追加評価は完全な出力検証と自生成変換対の比較という対象・条件で位置付ける．

\subsection{同一組立条件によるSWAMPED変換対の評価}
\label{sec:swamped-controlled}
SWAMPEDの公開変換コードの固定コミット\texttt{04e2bcc}を用いて，確認済みの入力標本にnop挿入と整数定数分割を適用した．これは今回生成した変換版であり，既報の変換済み配布標本を再評価したものとは区別する．変換関数の本体を保持し，WAT命令列を渡すアダプタを用いた．位置選択は元コードと同じbeta(1,1)分布を用い，再現用に乱数生成器を固定した．選択比は関数ごとの候補命令の0.1とし，候補がある場合の最低1か所という元実装の規則を保持した．整数定数分割では丸め誤差を混入させないため，浮動定数を含むvariant\_2関数を適用範囲から除いた．

nopはスタック・メモリ・制御を変えない命令であり，定数分割は整数定数$c$を$a+b=c$となる2定数と加算に置き換える．例えばvariant\_0関数の\texttt{i64.const 0}は，\texttt{i64.const -1}，\texttt{i64.const 1}，\texttt{i64.add}へ変わる．関数内の変更位置を保存し，この同値関係と基準版との対応を確認した．nopは30関数，定数分割は29関数に適用された．ここでは計測用hookやカウンタ更新を挿入していない．両手法の構造変更は保証されるが，解析の困難さや検出回避の効果を示す実験ではない．

WABT 1.0.42で基準と変換版を同じ方法で組み立てた．基準版のSHA-256は元配布物と完全一致した．3標本とも公表v0入力5件が一致し，importとexportの一覧も変わらなかった．同じホスト処理・メモリ配置・入力・1 workerで，各3秒のウォームアップと30秒以上の3反復を行った．各反復で実行順をずらし，初期化に使う時刻を固定した．全9ブロックで最終出力と，測定後の公表入力5件への出力が一致した．外部通信を無効にしており，プール提出も採掘報酬も発生しない．
\begin{table}[t]
\centering\small
\caption{SWAMPED入力への自生成変換の同条件比較（各30秒以上，3反復）}
\label{tab:swamped-controlled}
\begin{tabular}{lrrr}
\hline
変換 & H/s中央値 & 最小--最大 & 命令数比 \\
\hline
基準版 & 32.39 & 27.09--32.73 & 1.000 \\
nop挿入 & 31.08 & 29.56--32.72 & 1.097 \\
整数定数分割 & 31.53 & 28.94--32.08 & 1.034 \\
\hline
\end{tabular}
\end{table}
\begin{figure*}[t]
\centering
\includegraphics[width=\textwidth]{../results/runtime-mining/cross-dataset/hash-structure-comparison.pdf}
\caption{各実装の基準版に対するサイズ・命令数・速度の比．左はxmr-wasmの配布物比較，右はSWAMPED入力への自生成変換．速度の棒は中央値同士の比，白丸は各反復の変換版の速度を同じ反復の基準版の速度で割った3点である．同じ反復内でも順に測定しており，同時測定や信頼区間ではない．最適化条件が異なるxmr-wasm偽制御フロー版は図から除き，表に残した．}
\label{fig:swamped-controlled}
\end{figure*}

図\ref{fig:swamped-controlled}は，静的な増加量と速度変化の関係を実装ごとに示す．xmr-wasmでは，命令置換版の命令数・サイズの増加が平坦化版より大きい一方，速度低下は平坦化版のほうが大きい．白丸で示す各反復の基準比でも，平坦化版と命令置換版は3回とも1未満であった．一方，分割設定版とSWAMPEDの両変換では1の上下に変わる．基準版の変動を含めて差の方向を確認するが，時間経過や条件差を除いた因果効果とは扱わない．両パネルの変換手法と実装は異なり，実装間の絶対速度や手法の優劣は比較しない．

nop版のサイズ比は1.042，命令数比は1.097，速度中央値比は0.960であった．整数定数分割版では，それぞれ1.024，1.034，0.973となった．一つの入力標本と各手法の一つの乱数設定による測定であり，この変化率を全標本へ一般化しない．追加命令の静的増加が，そのまま同じ割合の性能低下になるとは判断できない．反復の最小・最大の範囲は3標本で重なり，この3反復から小さな中央値差を確定的な性能低下とは扱わない．また，挿入nopを除去し，分割式を剰余演算で元の定数へ戻すと，モジュール全体のWAT命令列が基準版と完全一致することを確認した．Wasmエンジンによる最適化の関与は考えられるが，生成された機械語を確認していないため，命令が除去されたとは断定しない．

測定順序も照合すると，基準版の速度は1回目の27.09 H/sから3回目の32.73 H/sへ20.8\%上がっていた．各反復内の基準版に対する観測比はnop版0.960--1.091，定数分割版0.973--1.068で，どちらも差の方向が入れ替わった．循環する測定順序でも時間経過の影響は完全には除けず，ウォームアップ，動作周波数，周辺負荷などの寄与は測定していない．少数反復の小さな中央値差を難読化の負荷と確定できない根拠として扱う．

\subsection{RQ2への回答}
\label{sec:rq2-answer}
評価した入力の範囲では，計算できた変換版も期待するハッシュ出力を保持した．一方，出力を照合できない配布標本もあり，実行制約と計算結果の不一致を区別した．速度と静的な増加量の関係は一様ではない．xmr-wasmの配布比較では平坦化と置換の増加量・速度低下の順序が一致せず，同一組立条件のSWAMPED自生成対でも，命令数比と速度比を別に確認する必要があった．各標本の変換範囲，最適化，ホスト処理を記録した範囲での結果であり，すべての採掘実装や難読化手法に共通する性能影響を確定したものではない．

\section{考察}
\subsection{仕組みと静的指標の関係}
今回の結果では，静的な増加量と実行時の負荷をつなぐために，変換された処理と測定入力が通る経路を照合する必要があった．例えば，xmr-wasmの置換版は平坦化版より命令数が増えたが，速度低下の順序は逆であった．さらに，公開hashラッパで確認した置換式やvariant 4以上に限定された処理は，今回のv0入力では通らない．従って，ファイル全体の命令数と変換箇所の存在だけで実行負荷を説明すると，測定していない処理の効果を混ぜる．これは今回の公開標本と条件に基づく結果であり，平坦化と置換の一般的な優劣ではない．

SWAMPEDの自生成対では，nop挿入と整数定数分割による命令数比が1.097と1.034であるのに対し，速度中央値比は0.960と0.973であった．反復範囲が重なり，反復内の基準比は1の上下に変わるため，小さな速度差を確定的な難読化効果としない．一方，xmr-wasmの配布比較では歴史的なビルド条件を統制できず，特に偽制御フロー設定版の最適化条件が異なる．実際の変換対で条件をそろえる評価と，配布物の記述的比較を区別することが，測定結果を再利用する際の条件になる．静的なツール間分布の違いも，性能や検出回避の順位を単独では与えない．

\subsection{複数のデータセットで確認できた範囲}
研究用に生成したデータでは，変換前後のファイルを使って手法を比較できる．一方，収集されたファイルからは公開時点の実装を調べられるが，変換の履歴が分からない場合がある．Cryptic Bytesでは設定横断の静的な変化と代表機構を，Proof-of-Theftでは配布設定と機構を照合した2つの実装の出力・速度を扱う．SWAMPEDでは入力標本の完全なハッシュ計算を確認し，公開変換コードによるnop挿入・整数定数分割の自生成対で機構，出力，速度を対応付けた．これらの実装間の独立性は未確定であり，データセット名の違いを独立した再現と数えない．収集されたファイルではコードの特徴の割合を調べ，別の範囲の情報を補っている．複数の資料に含まれるすべてのファイルで，すべての手法を確認したわけではない．

一般用途の資料や計測用コードの例は，測定結果を正しく理解するために必要だが，研究の主題ではない．本稿では分析方法の補助資料として位置付け，RQ1では手法の種類と収録内容，RQ2ではコード変化と正しさを確認したハッシュ計算性能の関係に焦点を合わせた．今後対象を広げる際も，加工前の採掘版，難読化版，計測版を区別する必要がある．

配布ソースの照合では，webminerpoolhashcnとxmr-wasmのOpenAES設定ファイルがバイト単位で一致し，AESラウンド処理\texttt{oaes\_encryption\_round}もコメント・空白を除いた命令の記述が一致した．同処理の引数検査，バイト置換，行・列の変換，鍵とのXORを照合した．これは共通部品の再利用を示すが，採掘処理全体や生成バイナリが同じことは示さない．二つの配布実装で測定できたことを，完全に独立した暗号実装での再現とは解釈しない．未実測の追加候補cryptonight-hashでも，OpenAES設定ヘッダやGroestlの表を共有し，webminerpoolhashcnの16バイトのコピー・XOR補助処理と，xmr-wasmの一部の加算補助処理がコメント・空白を除いた記述で一致した．この候補を追加しても，共通部品を持たない実装での再現とは数えられない．一方，補助処理の共有だけで採掘処理全体が同じとも判断しない．

\section{限界と再現性}
本研究は入手できた公開ファイルを対象とし，すべてのデータセットを調べたとは主張しない．新たなWebサイトの巡回や，2025年10月から2026年10月までの変化を追う調査は行っていない．MineSweeperの過去のcrawled-sites配布物に含まれる36種類は，元のファイルのままでは検証に失敗した．そのため，特徴の割合を求める際は別のminer samplesを用いた．解析に失敗したことや，入手した範囲による偏りが残る．

難読化の仕組みを確認したのは，目的に合わせて選んだcn-0の一部，xmr-wasmの公開hashとその呼出し先の一部，SWAMPED入力に自生成した変換対，未実測のcryptonight-hash内部候補の一部である．すべての設定や変換版について分類の精度を評価したわけではなく，収集されたファイルで難読化を見落とす割合も評価していない．また，一般用途のプログラムの事例を，採掘プログラムで確認した件数には含めていない．解析と執筆にはCodexの支援を用いた．状態を更新した後の戻り先，仮想命令の読出し，算術式の導出，関数に渡す引数の対応は，人間による独立したレビューが必要である．

静的比較の元となるCryptic Bytesの10本は，互いに関連するCryptoNightの派生プログラムである．追加のローカル実測13バイナリもcn/0またはcn-picoに限られ，データセットをまたいで別の実装系統が含まれることは未確定である．他の採掘アルゴリズムでも同じ結果になるとは限らない．同じ条件の基準版と比較したが，ビルド条件の違いや前処理の影響を完全に分けた実験ではない．Proof-of-Theftの実行比較は役割とホスト処理を照合した部分集合に限られる．すべての補助関数がどのように作られたかや，名前の変換による効果も，今回のデータセットでは確定していない．

実行評価は一部標本に限られる．プールで受理された13件に加え公表テストベクトルも照合したが，対象入力は有限であり，全入力の正しさを保証しない．プール受理とローカル出力一致，ブラウザ測定とNode.js測定を区別する．xmr-wasmへの拡張と30秒以上の3反復は完了したが，長時間の追加反復，実装間の独立性の確認，元配布の難読化版をビルド条件から再現した対照は未完了である．SWAMPEDの自生成対では同じ組立条件の比較を追加した．命令数だけから性能を推定した数値は示さず，検出回避能力も測定していない．

結果を再現できるように，入力ファイルの一覧，SHA-256，取得元のリビジョン，配布条件，設定名，対応する基準版，静的解析で測った値，元のプログラムごとの集計，代表的な仕組みを確認した命令の位置を保存した．図は日本語フォントを埋め込んだベクター形式とし，元のCSVと作成用コードを残す．主要な主張と根拠の対応表，研究上の制約，投稿前の確認事項を付属資料として整理する．外部公開と投稿は本作業に含めない．大学PCでの新たな実行は指導教員の許可確認後とし，許可された場合も報酬を発生させないローカル評価に限定する．既存のプール受理記録と追加の公表ベクトル照合を区別する．

emcc-obfの公開DockerfileはEmscriptenとBinaryenをコミットで固定する一方，Hikariをブランチから取得し，サブモジュールをリモート更新する\cite{emccobfdocs}．参照先のHikariリポジトリも本調査時点では取得できず，この手順だけでは配布物を生成した変換コードの版を特定できない．別のforkで再生成できても，元環境の再現と同一視せず条件変更として扱う必要がある．

\section{まとめ}
公開されたクリプトジャッキング関連のWasm資料を調べ，研究用に使われた変換設定と，採掘プログラムで確認できる難読化の仕組みを整理した．RQ1では，代表的な箇所で，平坦化，仮想命令を読みながらの処理，算術式の置換を確認した．一方，間接呼出しは通常の暗号処理にもあり，変換設定版で観測した呼出し先の取得方法だけでは，難読化による追加と断定できない箇所を区別した．資料によって含まれる設定と確認できた根拠の範囲が異なり，収集されたファイルで各手法が使われている割合は確定していない．

RQ2では，基準版を含む20,460件の比較を，元のプログラムごとにまとめた．命令数，ファイルサイズ，コードの構造，関数呼出し，メモリを読み書きする命令の変化は，手法と配布条件によって異なった．仮想化ではdataセクションの増加に加え，比較の基準となるファイルの大きさもサイズ比に影響した．これらは実行時の性能評価とは区別される結果である．本研究では，後続の評価で再利用できる変換前後の対応と，ファイルを実行せずに確認した根拠を示した．実行評価では一部の通常版と変換版でshare受理を記録し，Proof-of-Theftの5標本で13件の出力一致と短時間の反復速度測定を行った．xmr-wasmの5標本でも公表された期待値との一致と30秒以上の3反復を確認した．平坦化の機構を持つ版は命令数1.166倍，速度比0.678で，命令増加の順位と速度低下の順位は一致しなかった．SWAMPEDの入力標本に公開変換コードから生成した同条件の対照では，nopと整数定数分割の命令数比は1.097と1.034，速度中央値比は0.960と0.973であった．反復範囲は重なり，小差を確定的な効果としない．結果は今回の入力・標本・条件に限られ，元配布の難読化版を再現する統一ビルドと長時間の安定性の確認は未完了である．
\bibliographystyle{ipsjunsrt}
\bibliography{references}
\end{document}
